\section{Introduction}
\label{sec:intro}

Wearable biosensors and bedside instruments now record physiological signals continuously, in hospital wards as well as in home care and ambulatory rehabilitation~\cite{rudsari2025digital,khan2024implementing}. The resulting infrastructure, the Internet of Medical Things (IoMT), produces data that cannot all be moved to a central site within the privacy, latency and bandwidth limits of a care provider~\cite{nguyen2022federated,rauniyar2024federated}. Three technologies address parts of this problem. Federated learning (FL) trains a shared model from parameter updates instead of records~\cite{mcmahan2017communication,kairouz2021advances}. A digital twin (DT) maintains a patient-specific virtual model that is synchronized with sensor data~\cite{corral2020digital,kamelboulos2021digital}. Edge computing keeps the twin and part of the computation close to the patient. Combined, they give an architecture in which the twin of each patient, hosted on an edge node, takes part in a server-coordinated FL protocol while the records remain where they were collected.

The client scheduler of such a system decides, in every round, which devices train and transmit. Three requirements bear on that decision. The first is energy. A wearable device runs on a small battery whose first duty is monitoring, so the energy spent on local training and on the uplink has to leave a reserve for sensing~\cite{yang2021energy,arouj2022towards}. The second is participation. A scheduler that ranks clients by cost alone returns to the same inexpensive devices, and devices with slow processors, weak channels or large local data sets are rarely selected, so their data contribute little to the model~\cite{cho2022towards,huang2021efficiency}. The third is integrity. A client may submit an update that is corrupted by a fault or crafted by an adversary, and an unprotected average passes that update into the global model~\cite{blanchard2017machine,yin2018byzantine}. The three requirements interact: the cheapest clients form a small stable subset, rotation admits expensive clients, and a filter applied after selection changes who has effectively participated.

\rtag{E1, R1.1, R2.1}\begin{revblock}Each requirement has been studied, and several pairs have been combined. Energy-aware selection and resource allocation are established for wireless FL~\cite{nishio2019client,tran2019federated,yang2021energy}. Long-term participation fairness has been enforced with virtual queues and Lyapunov optimization~\cite{huang2021efficiency}, and FairFedCS combines such a queue with a contribution-based reputation in one index~\cite{shi2023fairfedcs}. Trust-FedGNN couples trust-weighted aggregation with energy-aware client scheduling~\cite{reis2026trustfedgnn}. Digital twins have been used to allocate resources in federated edge networks~\cite{lu2021communication,zhou2025energy}. The building blocks used in this article (energy prediction, stochastic quantization, the Jain index, the geometric median, cosine similarity and exponential reputation) are likewise well established. What remains open, and what this article addresses, is how energy, participation and reputation should enter \emph{one} scheduling decision for battery-constrained IoMT clients, what such a rule can be shown to guarantee, and how it compares with additive, sequential and drift-plus-penalty alternatives when everything except the rule is held fixed.\end{revblock}

This article presents ATEF-DT, an adaptive trust-aware and energy-fair framework for federated digital-twin healthcare. In each round a twin layer reports the device state and a forecast of the monitoring energy the device will need, a feasibility screen removes clients that cannot afford the round, and a single score ranks the remaining clients by predicted energy, participation count and reputation. Selected clients train locally, quantize their update with a bit-width assigned by the server and upload it. The server compares the received updates with a robust reference, updates the reputations, and aggregates the accepted updates with reputation weights after norm clipping. The contributions of this article are as follows.

\begin{itemize}
	\item \textbf{Patient-centric digital twin.} A Kalman twin tracks the physiological state of each patient and forecasts the monitoring energy over a short horizon; a feasibility screen admits a device to training only if this reserve remains intact.
	\item \textbf{Energy-aware client selection.} A single weighted-product score ranks the feasible clients by predicted energy, participation and reputation. \rtag{E2, R1.2, R2.2}\rev{The score is a heuristic, not the solution of the long-run problem; it is proved to solve a per-round surrogate exactly, to select no dominated client and to be invariant to a rescaling of the criteria}.
	\item \textbf{Adaptive quantization and credibility-based aggregation.} The bit-width of each client follows the progress of training, its channel and its battery. The server scores the received quantized updates against a geometric-median reference \rtag{R1.4}\rev{and aggregates them with reputation weights that sum to one over the accepted set}.
	\item \textbf{Convergence of the applied update.} \rtag{E3, R1.3, R2.3}\rev{A non-convex bound is derived in which the participation bias depends on the distance between the effective aggregation weights and the data proportions, and vanishes exactly when the two coincide}.
	\item \textbf{Evaluation with an explicit scope.} ATEF-DT is compared with six reference methods and four multi-objective schedulers on real clinical data. \rtag{R1.5}\rev{Energy is reported against full participation and uniform sampling, downlink traffic is counted, accuracy is broken down by device tier, and the attack regimes in which the method fails are reported.}
\end{itemize}

\rtag{E4, R1.5, R2.5}\rev{The scope of the evidence is stated at the outset. The learning task uses one real data set, the Wisconsin Diagnostic Breast Cancer (WDBC) data. The devices, batteries, wireless channel and energy consumption are simulated from stated models, and the physiological streams used to exercise the twin layer are synthetic. Section~\ref{sec:experiments} lists what is real and what is simulated, and Section~\ref{subsec:limitations} discusses the consequences.}

The remainder of this article is organized as follows. Section~\ref{sec:related} reviews related work. Section~\ref{sec:model} presents the system model. Section~\ref{sec:convergence} gives the convergence analysis. Section~\ref{sec:experiments} describes the experimental evaluation, Section~\ref{sec:results} presents the results and discusses their limitations, and Section~\ref{sec:conclusion} concludes.
